\documentclass{article}

\PassOptionsToPackage{numbers}{natbib}
\usepackage[preprint]{neurips_2026}

\usepackage[utf8]{inputenc}
\usepackage[T1]{fontenc}
\usepackage{hyperref}
\usepackage{url}
\usepackage{booktabs}
\usepackage{amsfonts}
\usepackage{amsmath}
\usepackage{amssymb}
\usepackage{nicefrac}
\usepackage{microtype}
\usepackage{xcolor}
\usepackage{graphicx}
\usepackage{enumitem}

\title{PIGNN-Attn-LS with K-Hop Gaussian Diffusion:\\
  Extending Edge-Aware Attention for Multi-Hop Power-Flow Context}

\author{%
  Changhun Kim
}

\begin{document}

\maketitle

%==============================================================================
\section{Background: Edge-Aware Self-Attention Solver (PIGNN-Attn-LS)}
%==============================================================================

\subsection{Unrolled Correction Steps}

The base model (PIGNN-Attn-LS, \texttt{GNSMsg\_SelfAttention\_armijo.py})
solves AC power flow by unrolling $K$ iterative correction steps over a graph
$\mathcal{G} = (\mathcal{V}, \mathcal{E})$ with $N$ buses.
At each step $k$, the power mismatches are computed from the current voltage
estimate $(v^{(k)}, \theta^{(k)})$:

\begin{equation}
  V_c = v \cdot e^{j\theta}, \quad
  I_c = Y V_c, \quad
  S_c = V_c \odot I_c^*, \quad
  \Delta P^{(k)} = P^{\mathrm{set}} - \mathrm{Re}\{S_c\}, \quad
  \Delta Q^{(k)} = Q^{\mathrm{set}} - \mathrm{Im}\{S_c\}
\end{equation}

with slack/PV masking applied ($\Delta P$ zeroed at the slack bus;
$\Delta Q$ zeroed at slack and PV buses).

\subsection{Physics-Informed Node Features}

The per-bus node state at step $k$ is:

\begin{equation}
  x_i^{(k)}
  =
  \bigl[
    v_i^{(k)},\;
    \theta_i^{(k)},\;
    \Delta P_i^{(k)},\;
    \Delta Q_i^{(k)},\;
    m_i^{(k)}
  \bigr]
  \in \mathbb{R}^{4+d}
\end{equation}

where $m_i^{(k)} \in \mathbb{R}^d$ is a learned latent memory vector.
The physics residuals $(\Delta P, \Delta Q)$ act as \emph{known-operator}
features: they inject the full nonlinear AC-PF equations directly into the
node representation at every step, so the GNN learns a residual-to-update
mapping rather than a blind function approximation.

\subsection{Edge-Aware Multi-Head Self-Attention}

Instead of uniform (isotropic) aggregation, the base model uses a
Transformer-style multi-head self-attention aggregator where each directed
edge $(j \to i)$ carries an \emph{edge bias} derived from line admittances.

For head $h$ with per-head dimension $d_h = d_{\mathrm{model}} / H$:

\begin{equation}
  q_i^{(h)} = W_Q^{(h)} x_i, \quad
  k_j^{(h)} = W_K^{(h)} x_j, \quad
  v_j^{(h)} = W_V^{(h)} x_j
\end{equation}

\begin{equation}
  \alpha_{ij}^{(h)}
  =
  \mathrm{softmax}_j
  \!\left(
    \frac{\langle q_i^{(h)}, k_j^{(h)} \rangle}{\sqrt{d_h}}
    + \beta_{ij}^{(h)}
  \right),
  \qquad
  \beta_{ij}^{(h)} = f_{\mathrm{edge}}(\ell_{ij})
\end{equation}

The edge bias $\beta_{ij}^{(h)}$ is a learned function of the
9-dimensional edge feature vector $\ell_{ij}$, which encodes the
physical line parameters:
$[y_{\mathrm{series},ij}^{\mathrm{re}},\; y_{\mathrm{series},ij}^{\mathrm{im}},\;
y_{\mathrm{shunt},i}^{\mathrm{re}},\; \ldots,\; \tau_{ij},\; \phi_{ij},\;
\mathbb{1}_{\mathrm{trafo}}]$.

This injects the admittance structure of the Jacobian $J(x)$ into attention
scores, making aggregation \emph{anisotropic} in a physically meaningful way:
lines with large admittances receive higher attention, mirroring the
Newton--Raphson Jacobian sensitivity structure.

The aggregated context is projected to produce voltage/angle corrections:

\begin{equation}
  \Delta\theta_i^{(k)} = W_\theta^{(k)} h_i^{(k)}, \quad
  \Delta v_i^{(k)} = W_v^{(k)} h_i^{(k)}, \quad
  \Delta m_i^{(k)} = \mathrm{LayerNorm}\!\bigl(\tanh(W_m^{(k)} h_i^{(k)})\bigr)
\end{equation}

An Armijo backtracking line search then selects the step size $\alpha$
that satisfies the sufficient-decrease condition on the power-mismatch
$\ell^\infty$-norm before accepting the update.

%==============================================================================
\section{Extension: K-Hop Gaussian Diffusion
  (\texttt{GNSMsg\_EdgeSelfAttnKHop})}
%==============================================================================

\subsection{Motivation}

The base attention model aggregates only \emph{direct} (1-hop) neighbors at
each message-passing step.  In power networks, a voltage disturbance
propagates electrically across multiple buses in a single physical instant;
the Jacobian $J(x)^{-1}$ therefore couples every bus to multi-hop
neighbors.  With only $K$ unrolled steps and 1-hop attention per step, the
effective receptive field is at most $K$ hops---which may be insufficient for
large grids without many iterations.

To address this, \texttt{GNSMsg\_EdgeSelfAttnKHop} augments the node state
with a \emph{K-Hop Gaussian (KHG) diffusion} \cite{zhang2026khg} of the
current node features before passing them to the attention blocks.  This
gives each bus an enriched view of its multi-hop electrical neighborhood
without increasing the depth of the attention network.

\subsection{Transition Matrix from the Admittance Matrix}

Rather than using a raw (binary) adjacency matrix, the model derives the
transition matrix directly from $Y$, the complex admittance matrix of the
grid.  Taking the element-wise absolute value and zeroing the diagonal:

\begin{equation}
  A_{ij} = |Y_{ij}|, \quad A_{ii} = 0
\end{equation}

$A_{ij}$ reflects the coupling strength between buses $i$ and $j$:
larger admittance means stronger electrical coupling.  Two options are
available via \texttt{--khop\_source}:
\begin{itemize}[noitemsep]
  \item \texttt{yabs} (default): $A_{ij} = |Y_{ij}|$, preserving
    admittance magnitudes as edge weights.
  \item \texttt{adj}: $A_{ij} = \mathbf{1}[|Y_{ij}| > 0]$, binary
    topology only.
\end{itemize}

The matrix $A$ is then normalized to a stable transition operator $T$
via \texttt{--khop\_norm}:
\begin{align}
  \text{row:} \quad & T = D_r^{-1} A, \quad D_r = \mathrm{diag}(A\mathbf{1})
  \label{eq:row}\\
  \text{sym:} \quad & T = D_r^{-1/2} A\, D_r^{-1/2}
  \label{eq:sym}\\
  \text{col:} \quad & T = A\, D_c^{-1}, \quad D_c = \mathrm{diag}(A^\top\mathbf{1})
  \label{eq:col}
\end{align}

Row normalization (default) makes each row sum to 1, so $T$ is a
row-stochastic matrix: it distributes a bus's information evenly weighted
by admittance to its neighbors.

\subsection{K-Hop Gaussian Diffusion Kernel}

Following \cite{zhang2026khg}, we apply $T$ iteratively to propagate
information across $K_{\mathrm{hop}}$ hops, weighting each hop $i$ by a
Gaussian that decays with distance:

\begin{equation}
  w(i, \sigma) = \exp\!\left(-\frac{i^2}{2\sigma^2}\right), \quad
  i = 1, \ldots, K_{\mathrm{hop}}
\end{equation}

where $\sigma > 0$ (\texttt{--khop\_sigma}) controls the effective range.
Small $\sigma$ keeps information local; large $\sigma$ extends it further.

The diffused feature matrix is the normalized weighted sum of hop-wise
propagations:

\begin{equation}
  D_K X
  =
  \frac{\displaystyle\sum_{i=1}^{K_{\mathrm{hop}}} w(i,\sigma)\; T^i X}
       {\displaystyle\sum_{i=1}^{K_{\mathrm{hop}}} w(i,\sigma)}
  \label{eq:dkx}
\end{equation}

where $T^i X$ is computed recursively as $H^{(i)} = T \cdot H^{(i-1)}$
with $H^{(0)} = X$.  The denominator $Z = \sum_i w(i,\sigma)$ normalizes
so the output stays on the same scale as the input regardless of
$K_{\mathrm{hop}}$ or $\sigma$.

Intuitively, $D_K x_i$ is a smoothed average of the node states in
the $K_{\mathrm{hop}}$-hop neighborhood of bus $i$, where closer buses
(stronger electrical coupling) contribute more.

\subsection{Integration into the Attention Solver}

The only change to the forward pass relative to the base model is how the
input to \texttt{in\_proj} is constructed.  In the base model:

\begin{equation}
  \text{(base)}\quad
  z_i^{(k)} = W_{\mathrm{in}}\, x_i^{(k)}
  \in \mathbb{R}^{d_{\mathrm{model}}}
\end{equation}

In the K-hop variant, the raw node state and its diffused counterpart are
\emph{concatenated} and projected together:

\begin{align}
  x_i^{(k)} &= \bigl[v_i, \theta_i, \Delta P_i, \Delta Q_i, m_i\bigr]
  \in \mathbb{R}^{4+d}
  \label{eq:nodestate}\\
  \tilde{x}_i^{(k)} &= \bigl(D_K X^{(k)}\bigr)_i
  \in \mathbb{R}^{4+d}
  \label{eq:diffused}\\
  z_i^{(k)} &= W_{\mathrm{in}}\,
  \bigl[x_i^{(k)} \;\|\; \tilde{x}_i^{(k)}\bigr]
  \in \mathbb{R}^{d_{\mathrm{model}}}
  \label{eq:proj}
\end{align}

where $\|\,$ denotes concatenation, so $W_{\mathrm{in}} \in
\mathbb{R}^{d_{\mathrm{model}} \times 2(4+d)}$.
Consequently, the input projection dimension doubles from $4+d$ to
$2(4+d)$, which is the \textbf{only} architectural change relative to the
base model.

After projection, the same \texttt{EdgeSelfAttnBlock} layers, output heads,
Armijo line search, and physics loss are used unchanged (they are directly
imported from \texttt{GNSMsg\_SelfAttention\_armijo}).

\subsection{What Each Component Contributes}

\begin{center}
\begin{tabular}{lll}
\toprule
Component & Where & Role \\
\midrule
$x_i^{(k)}$ & local (1-hop) & current bus state + physics residuals \\
$\tilde{x}_i^{(k)} = D_K x_i^{(k)}$ & multi-hop & smoothed neighborhood context \\
Edge bias $\beta_{ij}$ & attention & admittance-weighted message routing \\
Armijo line search & update & sufficient-decrease enforcement \\
\bottomrule
\end{tabular}
\end{center}

The diffusion $D_K$ acts as a \emph{preprocessing} module
\cite{zhang2026khg}: it is computed once per iteration step before the
attention blocks, not learned end-to-end.  This means the attention
mechanism still handles the adaptive, physics-aware aggregation, but each
bus now enters attention already aware of what is happening several hops
away---an electrical version of ``pre-smoothing.''

\subsection{Hyperparameters}

\begin{center}
\begin{tabular}{lll}
\toprule
Argument & Default & Meaning \\
\midrule
\texttt{--khop\_K} & 3 & maximum hop distance $K_{\mathrm{hop}}$ \\
\texttt{--khop\_sigma} & 1.5 & Gaussian decay scale $\sigma$ \\
\texttt{--khop\_norm} & \texttt{row} & transition matrix normalization \\
\texttt{--khop\_source} & \texttt{yabs} & edge weights from $|Y|$ or binary \\
\bottomrule
\end{tabular}
\end{center}

Setting \texttt{--khop\_K 0} recovers the base model exactly (the diffusion
returns a zero tensor, contributing nothing after projection).

%==============================================================================
\section{Summary of Differences from the Base Model}
%==============================================================================

\begin{center}
\begin{tabular}{lll}
\toprule
Aspect & PIGNN-Attn-LS (base) & K-Hop variant \\
\midrule
Node input to \texttt{in\_proj} & $x_i \in \mathbb{R}^{4+d}$ &
  $[x_i \| D_K x_i] \in \mathbb{R}^{2(4+d)}$ \\
\texttt{in\_proj} width & $4+d$ & $2(4+d)$ \\
Receptive field per step & 1-hop (attn.) &
  $K_{\mathrm{hop}}$-hop (diff.)\ +\ 1-hop (attn.) \\
Transition matrix source & --- & $|Y|$ (admittance-weighted) \\
Attention blocks & imported \texttt{EdgeSelfAttnBlock} &
  same, unchanged \\
Physics loss & same & same \\
Armijo line search & same & same \\
\bottomrule
\end{tabular}
\end{center}

\begin{thebibliography}{9}

\bibitem{zhang2026khg}
X.~Zhang, P.~Wang, D.~Li, A.~Huang, Z.~Chen, and Y.~Yang,
``Enhanced graph neural networks using K-hop Gaussian diffusion,''
in \emph{ICASSP 2026}, 2026.

\bibitem{kim2026pignn}
C.~Kim et al.,
``Physics-informed GNN for medium-high voltage AC power flow with
edge-aware attention and line search correction operator,''
in \emph{ICASSP 2026}, 2026.

\end{thebibliography}

\end{document}
